\chapter{Literature Review}\label{ch:A}

\section{Photoelectrochemical Water Splitting: Fundamentals}

\subsection{Semiconductor Photoelectrodes}

A PEC cell consists of a semiconductor photoelectrode in contact with an aqueous
electrolyte, typically with a Pt counter electrode.
When illuminated, the semiconductor absorbs photons with energy $E \geq E_g$
(bandgap energy), creating electron--hole pairs.
Band bending at the semiconductor--electrolyte interface separates the charges:
holes migrate to the surface for OER, while electrons flow through the external
circuit to the cathode for HER~\cite{gratzel2001photoelectrochemical}.

The ideal semiconductor for unassisted water splitting requires:
\begin{itemize}
  \item A bandgap between 1.8 and 2.2\,eV to absorb a significant fraction of the
        solar spectrum while providing sufficient overpotential for both
        half-reactions~\cite{bolton1978limiting}.
  \item Conduction and valence band positions that straddle the water
        oxidation ($+1.23\,\text{V}$ vs.\ RHE) and reduction (0\,V vs.\ RHE)
        potentials.
  \item Chemical stability in the electrolyte under prolonged illumination.
  \item Efficient charge transport and low surface recombination.
\end{itemize}

\subsection{Key Performance Metrics}

Three metrics are universally reported in PEC literature and constitute the
prediction targets in this work:

\paragraph{Photocurrent Density $J_{ph}$} (mA/cm\textsuperscript{2}) is measured
by linear sweep voltammetry under AM~1.5G illumination.
It reflects the rate of charge transfer to the electrolyte and serves as the
primary proxy for solar absorption and charge collection efficiency.
State-of-the-art values exceed 6\,mA/cm\textsuperscript{2} for BiVO\textsubscript{4}
\cite{kim2014bismuth} and reach $>$20\,mA/cm\textsuperscript{2} for GaN nanowires
\cite{kibria2015multibandgap}.

\paragraph{Solar-to-Hydrogen Efficiency (STH)} is the gold standard for overall
device efficiency (Equation~\ref{eq:sth}).
The theoretical maximum for a single-junction device with $E_g = 2.0$\,eV is
approximately 16.8\% \cite{walter2010solar}.
The highest experimentally reported STH for a single photoelectrode under
AM~1.5G without applied bias is around 12\% for III-nitride nanowires.

\paragraph{Hydrogen Evolution Rate (HER)} ($\mu$mol\,h\textsuperscript{-1}\,cm\textsuperscript{-2})
is measured by gas chromatography.
It is related to photocurrent through Faraday's law:
\begin{equation}
  \text{HER} = \frac{J_{ph} \times 3600}{2 \times F} \times \eta_F \times 10^6
  \label{eq:her}
\end{equation}
where $F = 96485$\,C/mol.
In practice, HER is reported less frequently than $J_{ph}$ because it requires
a gas-tight cell and calibrated gas chromatograph.

\subsection{Principal Material Systems}

Table~\ref{tab:materials} summarises the main photoelectrode materials covered
in this study.

\begin{table}[h]
\centering
\caption{Common PEC photoelectrode materials and typical performance ranges.}
\label{tab:materials}
\begin{tabular}{llcc}
\toprule
Material & Bandgap (eV) & Typical $J_{ph}$ (mA/cm\textsuperscript{2}) & Role \\
\midrule
BiVO\textsubscript{4}       & 2.4--2.5 & 1--7   & Photoanode \\
$\alpha$-Fe\textsubscript{2}O\textsubscript{3} & 2.0--2.2 & 0.3--4 & Photoanode \\
WO\textsubscript{3}         & 2.6--2.8 & 1--6   & Photoanode \\
TiO\textsubscript{2}        & 3.0--3.2 & 0.2--2 & Photoanode \\
ZnO                         & 3.2--3.4 & 0.5--3 & Photoanode \\
Cu\textsubscript{2}O        & 1.9--2.1 & 1--5   & Photocathode \\
GaN / InGaN                 & 0.7--3.4 & 5--22  & Photoanode \\
Ta\textsubscript{3}N\textsubscript{5} & 2.0--2.1 & 0.3--4 & Photoanode \\
Si                          & 1.1     & 8--30  & Photocathode \\
\bottomrule
\end{tabular}
\end{table}

\section{Machine Learning for Materials Discovery}

\subsection{General Approaches}

Materials informatics applies statistical and ML methods to accelerate
the discovery of new functional materials.
Pilania et al.~\cite{pilania2013accelerating} demonstrated that random
forest models trained on DFT-computed features could predict formation
energies with mean absolute errors below 0.1\,eV/atom.
Isayev et al.~\cite{isayev2017universal} introduced universal fragment
descriptors for predicting band structures across thousands of inorganic
compounds.

In photocatalysis, Ye et al.~\cite{ye2018deep} applied convolutional neural
networks to crystal structure images to predict photocatalytic activity.
Schmidt et al.~\cite{schmidt2019recent} reviewed supervised regression models
for bandgap prediction from composition alone, reporting R\textsuperscript{2}
values of 0.85--0.94 on large DFT datasets.

\subsection{ML Applied to PEC Systems}

Compared with solid-state properties, PEC performance prediction is harder
because the output depends not only on the material but also on electrolyte
chemistry, electrode geometry, and illumination conditions.
Wu et al.~\cite{wu2019machine} trained a gradient-boosted model on 150
published OER catalyst results, showing that material composition alone
could predict overpotential with an MAE of 27\,mV.

For water splitting specifically, Chen et al.~\cite{chen2019graph} applied
graph neural networks to predict the photocatalytic hydrogen evolution
rate of organic semiconductors from molecular structure.
However, most existing studies focus on either purely computational
descriptors (DFT band structures, formation energies) or organic
photocatalysts, leaving the experimental PEC inorganic literature
largely untapped.

\subsection{Ensemble Tree Methods}

Gradient boosting and random forests are the dominant methods in tabular
data competitions and materials property prediction for datasets of this
scale ($n < 10^4$) \cite{grinsztajn2022tree}.

\textbf{XGBoost}~\cite{chen2016xgboost} builds an additive ensemble of decision
trees by sequentially fitting residuals with a regularised objective:
\begin{equation}
  \mathcal{L}(\phi) = \sum_i l(y_i, \hat{y}_i) + \sum_k \Omega(f_k)
\end{equation}
where $l$ is a differentiable loss, $\hat{y}_i$ is the predicted value,
and $\Omega(f_k) = \gamma T + \frac{1}{2}\lambda \|w\|^2$ penalises tree
complexity.

\textbf{Random Forest}~\cite{breiman2001random} averages predictions of $B$
de-correlated trees, each trained on a bootstrap sample with feature
subsampling at each split.
Beyond prediction, the variance across individual trees provides a natural
uncertainty proxy \cite{wager2014confidence}.

\section{Feature Engineering for PEC}

Physics-informed feature engineering transforms domain knowledge into
numerical features that improve model interpretability and data efficiency.
For PEC systems, the most informative features encode electrochemical
driving force, light absorption, and surface kinetics.
Huo et al.~\cite{huo2019unified} showed that including thermodynamic
descriptors (bandgap, electron affinity) alongside synthesis conditions
improved prediction of heterogeneous photocatalytic activity by 30\%
compared to composition-only models.

The practice of including interaction features (e.g., $E_g \times \text{pH}$)
is supported by the Butler--Volmer model:
the flat-band potential of a metal oxide shifts with pH at $-59$\,mV/pH unit
(Nernst equation), meaning the net driving force depends on the
interplay between bandgap and solution acidity.

\section{Gap Analysis}

The following gaps motivate the present work:
\begin{enumerate}
  \item No publicly available, unified ML-ready dataset covers a broad
        range of inorganic PEC photoelectrode materials with experimental
        performance metrics.
  \item Existing PEC ML studies do not address data sparsity in STH and HER
        targets~— the two most practically relevant metrics.
  \item Deployed web applications that expose PEC prediction models
        interactively have not been reported in the literature.
  \item Uncertainty quantification in PEC property prediction (model
        confidence, variance across estimators) is absent from prior work.
\end{enumerate}
